% STATUS: written 2026-11-04 (Chuang--Lee--Liu--Tsai 2010 and Chuang 2017 read in full)
% Sources read in full: Chuang, Lee, Liu and Tsai 2010 (pp. 157--178); Chuang 2017 (pp. 2837--2847).
% The classical theory (Kharchenko 1978) is recalled in Section~\ref{sec:prime-semiprime:kharchenko}; this section reports the generalized and skew versions only.
\section{Generalized and skew differential identities}\label{sec:generalized-twisted:differential-identities}

A differential identity is an expression involving unknowns, ring elements and
derivations as unary operations \cite[p.~155]{Kharchenko1978}. Kharchenko built
the theory of these identities for prime rings, and
Section~\ref{sec:prime-semiprime:kharchenko} reports it. The present section
reports two later extensions \cite[p.~2838]{Chuang2017Auto}. The first has
derivations that act through the twist of a generalized derivation
\cite[p.~885]{DeFilippisFosner2012}. The second has a twist that is a skew
derivation of the Ore extension type \cite[p.~2838]{Chuang2017Auto}.

The first step is to translate an identity about a derivation into an identity
about an Ore extension \cite[p.~172]{ChuangLee2010Higher}. The translation rests
on the commutator formula for the variable $x$
\cite[p.~172]{ChuangLee2010Higher}.

\begin{remark}[{\cite[p.~172]{ChuangLee2010Higher}}]\label{rem:gen:identity-translate}
In $Q[x;\delta]$ the derivation acts as the commutator with $x$. One has
$\delta(r)=[x,r]$ and, in general, % claim-ok: identity displayed from the cited page
\begin{equation*}
\delta^{i}(r)=\sum_{s=0}^{i}(-1)^{s}\binom{i}{s}x^{i-s}r\,x^{s}.
\end{equation*}
A differential identity $\sum_{ij}a_{ij}\delta^{i}(r)b_{ij}=0$ for all $r\in R$
is thereby an equality $\sum_{i}f_{i}(x)\,r\,g_{i}(x)=0$ with $f_{i},g_{i}\in
Q[x;\delta]$.
\end{remark}

The translation turns a question about the values of derivations into a
question about an Ore extension \cite[p.~172]{ChuangLee2010Higher}. The answer
is then a linear independence statement
\cite[Theorem~11, p.~172]{ChuangLee2010Higher}.

\begin{theorem}[{\cite[Theorem~11, p.~172]{ChuangLee2010Higher}}]\label{thm:gen:identity-independence}
For $f_{i},g_{i}\in Q[x;\delta]$, one has $\sum_{i=1}^{n}f_{i}\,r\,g_{i}=0$ for
all $r\in R$ if and only if $\sum_{i=1}^{n}f_{i}\otimes g_{i}=0$ in
$Q[x;\delta]\otimes_{C_{S}(R)}Q[x;\delta]$.
\end{theorem}

The theorem is the single-derivation case of a general statement about
$(R,R)$-bimodule endomorphisms \cite[p.~164]{ChuangLee2010Higher}. Chuang, Lee,
Liu and Tsai prove it through these endomorphisms
\cite[Theorem~3, p.~165]{ChuangLee2010Higher}. Every
$(R,R)$-bimodule endomorphism of the Ore extension is a series in the divided
powers (Theorem~\ref{thm:gen:ore-endomorphism}).

The sources also state the structural consequence that the same authors draw
for the module of the extension \cite[p.~169]{ChuangLee2010Higher}.

\begin{theorem}[{\cite[Corollary~5, p.~169]{ChuangLee2010Higher}}]\label{thm:gen:ore-quasi-injective}
The $(R,R)$-bimodule $S=Q[X;D]$ is quasi-injective in the sense that any
$(R,R)$-bimodule map from an $(R,R)$-bisubmodule into $S$ can be extended to an
$(R,R)$-bimodule endomorphism of $S$.
\end{theorem}

The skew case starts one step earlier, at the definition of the identity
itself \cite[p.~2838]{Chuang2017Auto}. The twist $\sigma$ acts on the
coefficients, and the Ore extension of
Definition~\ref{def:gen:ore-extension} carries both maps.

The two maps are not independent when one of them is inner
\cite[p.~2841]{Chuang2017Auto}. This is the mechanism by which a skew identity
reduces to an ordinary one \cite[p.~2841]{Chuang2017Auto}.

\begin{remark}[{\cite[p.~2841]{Chuang2017Auto}}]\label{rem:gen:skew-reduce}
Let $\delta$ be a $\sigma$-derivation with $\sigma=I_{u}$ for a unit $u$ of
$R$. Then $d:=u^{-1}\delta$ is an ordinary derivation, and $R[t;d]$ is
isomorphic to $R[t;\sigma,\delta]$ by the $R$-fixing map $t\mapsto u^{-1}t$. If
$\delta$ is the inner $\sigma$-derivation defined by $a$, so that
$\delta(r)=ar-\sigma(r)a$, then $(t-a)r=\sigma(r)(t-a)$, and
$R[t;\sigma,\delta]$ is isomorphic to the Ore extension of automorphism type
$R[t;\sigma]$ by $t\mapsto t+a$.
\end{remark}

The reduction shows where the twisted identities are genuinely new
\cite[p.~2838]{Chuang2017Auto}. Some skew derivations are inner, and the
reduction above removes their twist
\cite[p.~2841]{Chuang2017Auto}. Some are of automorphism type, and a
change of variable removes the other part. In the remaining case the skew
derivation keeps a twist \cite[p.~2841]{Chuang2017Auto}. No change of
variable removes it; its identities are identities of the twisted polynomial
ring \cite[p.~2841]{Chuang2017Auto}.

The two sections of this chapter meet at the same phenomenon
\cite[p.~172]{ChuangLee2010Higher}. The Ore extension is the coefficient
ring of the identity, and the identities of
Theorem~\ref{thm:gen:identity-independence} are the identities of that ring.
Section~\ref{sec:generalized-twisted:semi-left-reverse} turns to the weaker
product rules, where the identity is not a twist of the Leibniz rule but a
different rule altogether.
